\documentclass{article}
\usepackage{amsmath}
\title{A correct branch report need not qualify its old motion certificate}
\date{}
\begin{document}
\maketitle

\paragraph{Task and admitted reference.}
Keep the selected KTH planar warm state, two authored hidden offset
branches, old-link and stopped-body separation requirements, and
norm-bounded double integrator of the preceding studies. This time
collapse the old vertex at selected earlier times along the very
same cubic prehistory. The failed body parks at its actual cutoff
pose. Hold the horizon at $T=1.11$ \emph{after} each fault, not at a
global clock deadline. Still require all survivors at rest with
a new link of range $1.54$. The sign is assumed known before each
cut, and remains correct. The recalculated policies also receive an
exact trusted time offset in the known prehistory, enough to derive
that cut's positions and velocities from the conditional model.
This does not supply a synchronization or timestamp protocol,
does not authorize command reuse on an arbitrary perturbation, and
does not carry old RF clearance into another regime. Pinned source
material is identified in \verb|SOURCE.tsv|.

\paragraph{Obstruction to reusing an old command.}
If a controller executes the exact old acceleration tape from a
different initial velocity, their terminal velocity difference in
a double integrator is exactly their initial difference. Knowing
the same sign or even the current position cannot make it stop.
For a fault $0.01$ units early, the additional remote position
differs from the original stop by only about $0.00002661$,
while the residual speed of replaying the same command would be
about $0.00531556$. This terminal failure does not assert that the
physical objective is impossible. Nor does the observation error
allowance of the earlier sign classifier become a tube guarantee.

\paragraph{A state-conditioned patch that retains the old direction.}
For displacement $d$, initial velocity $v$, and a first pulse
lasting $s$ out of $T$, choose
\[
 a={2d\over sT}-{T+s\over sT}v,qquad
 b=-{2d\over (T-s)T}+{s\over (T-s)T}v.
\]
The resulting terminal velocity is zero and the displacement is
exactly $d$, coordinatewise. The two integration identities and the
common-impulse counterexample are checked in \verb|Two_Pulse.lean|.
Parameterize $d$ as a scalar times the original branch's post-brake
bridge direction, with opposite senses on the two components.
Starting position determines the smallest collinear total shift
that puts the selected bridge inside the range sphere. Fixed
transverse offset is charged by the sphere equation rather than
rounded to zero. Each input norm bound is a scalar interval after
separating velocity parallel and perpendicular to the direction.
Intersect the per-robot intervals within each component, allow a
different pulse duration per component, and choose two admissible
terminal shifts with the required sum. Intermediate formation
deformations are allowed but checked. Between all actual acceleration
switches, a quadratic chord enclosure checks every old established
edge and every clearance throughout the motion. Affine velocity
between switch cuts makes the endpoint speed cap sufficient.
The duration search is a finite constructive grid, not a claim of
exhaustiveness over all controls or over continuous fault epochs.

\paragraph{Which conclusions transport.}
The 20 tested cases use both branches and ten specified fault
advances. At advances 0.008 and 0.01 the earlier, recomputed
brake-then-translate witness and the equal-half two-pulse attempt
both miss the requirement. Retaining the old direction but
adjusting pulse durations produces a checked witness in all four
of those cases. A direct method that uses the current bridge axis
and the same pulse family has the same acceptance pattern on this
test, with no resource-complexity advantage claimed for either.
At advance 0.016 neither checked pulse method certifies service,
but the continuous full-information lower bound is still below
the deadline; physical feasibility there is unresolved. At
0.032 and 0.064 the signed projected displacement bound rules
out all full-information controllers under the same limits.
These distinctions are part of the result, not failures silently
removed from the family.

The experiment justifies reusing an axis and endpoint check,
not an old state/input witness, when the state underlying a valid
branch report changes. Extra sensing resources, uncertainty in
right activation, and live reliable timestamps are outside this
reference. The exact crash-controller equations of the named R6
source, an independent mission, and an R6 acceptance claim remain
open.
\end{document}
